\scanpage{6}\mspage{\unsure{\phantom{5}|5}}
\[
\underline{\mathrm{Hom}}_{\otimes\mathrm{AUC}}(S^{-1}C,G)\handwritten{\longrightarrow}\underline{\mathrm{Hom}}_{\otimes\mathrm{AUC}}^{S}(C,G)
\]
(where the superscript $S$ in the second member denotes the full subcategory formed by the AUC $\otimes$-functors $F:C\handwritten{\longrightarrow}G$ such that $\deleted{P}F(X)$ is invertib\unsure{\textperiodcentered|-}le for every $X\handwritten{\in}\mathrm{Ob}\ S$) is an \uline{isomorphism} \uline{of categories}.

We do not spell out here the construction, fairly natural, but whose step-by-step spelling out, \deleted{\rlap{xx}and}\inserted{with} the verification of all the required compatibilities, is very long and tedious. Let us only point out here that when $S$ is the set of all objects of $C$, $S^{-1}C$ is itself a Picard category (not necessarily strict)\unsure{\textsuperscript{.}|'}, which one may call the ``enveloping Picard category'' of $C$. This is probably the most important case. The invariants $\handwritten{\pi}_{0}$ and $\handwritten{\pi}_{1}$ \deleted{\unsure{the|of}}of this enveloping category deserve the notation $K^{0}_{\deleted{0}}(C)$ and $K^{\unsure{1|\deleted{0}1|\deleted{1}1}}_{\deleted{1}}(C)$ respectively. For example, when $C$ is the category of finitely generated projective modules over a ring $A$, \inserted{with the sum operation,} one can prove that these groups are none other than the classical invariants $K^{0}(A)$ and $K^{1}(A)$ studied by Grothendieck and by Whitehead-Bass. The same result is probably valid for every additive category $C$ (equipped with \deleted{\rlap{xx}the} the sum operation), This gives a remarkable interpretation of the invariants $K^{0}$ and $K^{1}$ in terms of a ``stabilized category'' \handwritten{(}of the additive category $C$\handwritten{)}, which is a Picard category. Let us only point out that, this stabilized category (even in the special case of finitely generated projective modules over the ring $A$) \deleted{\unsure{is}}\handwritten{be}ing practically never a \uline{strict} Picard category, these considerations seem to make it unlikely that one will manage to construct a canonical object in the derived category of the category (Ab), which is a chain complex satisfying $H_{0}(K_{\bullet})=K^{0}$, $H_{1}(K_{\bullet})=K^{1}$ - since such a construction (thanks to Deligne's theory) would provide us with precisely a canonical \uline{strict} Picard category whose invariants would be $K^{0}$ and $K^{1}$.

One can hardly escape the question of an analogous geometric interpretation of the higher invariants $K^{i}$. There is every reason to believe that such an interpretation is to\unsure{\textasciidieresis|''} be sought in the theory of Picard $n$-categories (in a suitable sense) and the construction of enveloping Picard $n$-categories of an arbitrary AUC $\otimes$-category (or of an AUC $\otimes$-$n$-category) $C$, whose ``homotopy'' invariants $\handwritten{\pi}_{i}$ would be the $K^{i}$ sought. Such developme\deleted{b}nts would obli-
% The live text ends here on this scan with the printed continuation hyphen.

\deleted{As \deleted{\unsure{in}}another open direction of research, suggested by the results sketched here, let us point out the construction of a $\otimes$-category of fractions of an AU $\otimes$-category (without a commutativity datum neces-}
% archive sheet number: 5
